\documentclass[11pt]{article}
\usepackage[margin=1in]{geometry}
\usepackage{amsmath,amssymb,amsthm}
\usepackage[hidelinks]{hyperref}
\newtheorem{theorem}{Theorem}
\newcommand{\CI}{C_{\mathrm I}}
\newcommand{\CII}{C_{\mathrm {II}}}
\title{A counterexample in every band of the parabola conjecture\\
\large A solution to MathDB 347624}
\author{Principia Math In Progress}
\date{18 September 2026}
\begin{document}
\maketitle
\begin{abstract}
Vielhaber's proposed quadratic lower bound for the inverse binary
question-mark function is false. At $y=7/8$ the inverse is $4/5$,
whereas the proposed lower bound is $13/16$. More generally, the
bound fails by exactly $1/(80\cdot2^k)$ at $y=7/2^{k+3}$, for
every integer $k\geq0$. All evaluations follow directly from the
defining finite binary codes.
\end{abstract}

\section*{Task 1: the original problem and its full solution}
Write $\phi=\overline{?}_V^{-1}$ for the real-valued inverse in
Definition 2(iii) of \cite{V}. Conjecture 22 of that source asserts
that, whenever $k\geq0$ and $2^{-k-1}\leq y\leq2^{-k}$,
\begin{equation}\label{bound}
 \phi(y)\geq Q_k(y),\qquad
 Q_k(y)=\frac43\,2^k\left(y-\frac34\,2^{-k}\right)^2
               +y-\frac1{12}\,2^{-k}.
\end{equation}
It also specifies equality at the two band endpoints and its
three-quarter point. We disprove the inequality itself at interior
points distinct from all three specified points.

\subsection*{The coding convention}
For a positive integer $b=2^\ell+\sum_{j=0}^{\ell-1}a_j2^j$,
where $a_j\in\{0,1\}$, the source defines
\[
 \CI(b)=0^\ell1\overline a_{\ell-1}\cdots\overline a_0,
 \qquad
 \CII(b)=1^\ell0a_{\ell-1}\cdots a_0,
 \qquad \overline a=1-a.
\]
An exponent on a bit denotes repetition; an exponent zero denotes
the empty word. If $r=[b_1,\ldots,b_{2m}]\in(0,1)$ is an even-length
continued fraction, in the convention
$[b_1,b_2]=1/(b_1+1/b_2)$, concatenate
\[
 w=\CI(b_1)\CII(b_2)\cdots\CI(b_{2m-1})\CII(b_{2m}).
\]
Remove all trailing zeroes and then the final one to obtain the
address $v$. The corresponding dyadic argument is $y=(0.v1)_2$,
and $\phi(y)=r$. Equivalently, $y=(0.w)_2$, since removing trailing
zeroes does not change the binary fraction. This is Definition 2
and the address identification in Appendix 2 of \cite{V}.

\subsection*{The first counterexample}
The even-length expansion $4/5=[1,4]$ has
\[
 \CI(1)\CII(4)=1\,11000=111000.
\]
Its address is $v=11$, so $\phi(7/8)=4/5$. Independently, Figure 4
of the source gives precisely this address, argument, and value.
But
\[
 Q_0(7/8)=\frac43\left(\frac18\right)^2+\frac78-\frac1{12}
         =\frac{13}{16}
         =\frac45+\frac1{80}.
\]
Thus \eqref{bound} is false.

\subsection*{An exact failure in every band}
\begin{theorem}
For every integer $k\geq0$, put $y_k=7/2^{k+3}$. Then
\[
 \phi(y_k)=\frac4{5\cdot2^k},\qquad
 Q_k(y_k)=\frac{13}{16\cdot2^k},\qquad
 Q_k(y_k)-\phi(y_k)=\frac1{80\cdot2^k}>0.
\]
\end{theorem}
\begin{proof}
The case $k=0$ was proved above. For $k=1$,
\[
 \frac25=[2,2],\qquad \CI(2)\CII(2)=011\,100,
 \qquad (0.011100)_2=\frac7{16}.
\]
Now let $k\geq2$ and set $b=5\cdot2^{k-2}-1$. Since
$2^k\leq b<2^{k+1}$, its code uses $\ell=k$. Moreover,
\[
 b-2^k=2^{k-2}-1.
\]
Written with exactly $k$ bits, this last integer is $00\,1^{k-2}$;
its bitwise complement is $11\,0^{k-2}$. Hence
\[
 \frac4{5\cdot2^k}=\frac1{b+1}=[b,1],\qquad
 \CI(b)\CII(1)=0^k1110^{k-2}\,0.
\]
The associated dyadic value is
$2^{-k-1}+2^{-k-2}+2^{-k-3}=7/2^{k+3}$, proving the inverse
evaluation. In every case the address is $0^k11$.

Since $y_k=(7/8)2^{-k}$, it lies strictly inside the required band.
Substitution gives
\[
 Q_k(y_k)=2^{-k}\left(\frac43\left(\frac18\right)^2
                  +\frac78-\frac1{12}\right)
         =\frac{13}{16\cdot2^k}.
\]
Subtracting $4/(5\cdot2^k)$ gives the stated positive gap.
\end{proof}

This resolves the conjecture negatively. The obstruction uses its
stated coding convention and occurs in every band, including
$k=4$ mentioned in the source. It does not depend on an endpoint
interpretation, a numerical approximation, or the existence of
the extension away from dyadic arguments.

\section*{Task 2: the least uniform vertical correction}
Determine the smallest $c\geq0$ for which
$\phi(y)\geq Q_k(y)-c2^{-k}$ throughout every band. The theorem
forces $c\geq1/80$. Obtain a matching upper bound or identify
larger obstructions and the points approaching the supremum.

\section*{Task 3: the limiting shape of small bands}
For $1/2\leq t\leq1$, study $F_k(t)=2^k\phi(2^{-k}t)$.
Determine its limiting profile and a uniform error estimate,
using the first partial denominator and the remaining code.
The exact identity $F_k(7/8)=4/5$ supplies a consistency check.

\section*{Task 4: the set where the quadratic bound fails}
For each $k$, describe the set of $t\in[1/2,1]$ satisfying
$F_k(t)<\frac43(t-3/4)^2+t-1/12$. First classify dyadic
arguments by their finite codes, then determine the interval
structure and measure of the full failure set.

\section*{Task 5: sharp multiplicative bounds}
Determine $\inf_{0<y<1}\phi(y)/y$ and
$\sup_{0<y<1}\phi(y)/y$, with their equality cases. Compare the
result with the separate linear-bound Conjecture 21 of \cite{V}.
The present counterexamples have the constant ratio $32/35$ and
do not disprove its proposed lower ratio $8/9$.

\paragraph{Verification and status.}
The accompanying standard-library Python verifier uses exact
rational arithmetic and separately implements coding and decoding.
Its finite checks supplement the proof for every $k$. Tasks 2--5
are proposed directions and receive no solution credit. The full
20-page primary source and the documented bounded prior-claim
search were checked by the authoring agent; there has been no
independent peer review.

\begin{thebibliography}{9}
\bibitem{V} M. Vielhaber,
\emph{V Tree---Continued Fraction Expansion, Stern--Brocot Tree,
Minkowski's $?(x)$ Function In Binary: Exponentially Faster},
arXiv:2008.08020v1 (2020), Definitions 1--2, Conjecture 22,
and Figure 4. \url{https://arxiv.org/abs/2008.08020}.
\end{thebibliography}
\end{document}
